%% toc: P3 40 40.9
\section{Proofs for Ehrenfeucht-game proofs of asymptotic results}
\label{sec:annex-p03-w03-p03-072-game-proofs}

Throughout, languages have equality and no function or constant symbols,
and structures have nonempty domains.  Compactness
(\cref{thm:fol-compactness}) is used for countable languages.  Assembling a
strategy for infinite structures from pointwise choices of answers uses the
Axiom of Choice, as in \cref{thm:p03-w03-p03-038-pointed-ef}.

\begin{annexitem}{annex:p03-w03-p03-072-game-theorem}{thm:p03-w03-p03-072-game-theorem}{The game theorem}
\begin{proof}
\emph{The game is the pointed game with no parameters.}  Let \(\tau\) be the
set of relation symbols of \(L\).  In \cref{def:p03-w03-p03-038-positions},
\(G_t(G,\varnothing;H,\varnothing)\) has \(t\) rounds, each consisting of a
Spoiler choice in either structure and a Duplicator answer in the other, with
the pair appended to the position; these are the moves of
\(\mathrm{EHR}(G,H;t)\).  Duplicator wins a play of
\(G_t(G,\varnothing;H,\varnothing)\) when the final position is a partial
isomorphism, that is, when for every atomic formula in \(\mathrm{At}_t\) the
two tuples agree.  \(\mathrm{At}_t\) consists of the formulas
\(x_i=x_j\) for \(i,j\leq t\) and \(R(x_{i_1},\ldots,x_{i_l})\) for every
\(l\)-ary \(R\in\tau\) and all indices \(i_1,\ldots,i_l\leq t\).  Agreement on
these is exactly Spencer's winning condition: \(x_i=x_j\) iff \(y_i=y_j\),
and \(R(x_{i_1},\ldots,x_{i_l})\) iff \(R(y_{i_1},\ldots,y_{i_l})\).
Strategies and winning strategies are defined in the same way in both
places.  So Duplicator wins \(\mathrm{EHR}(G,H;t)\) iff she has a winning
strategy in \(G_t(G,\varnothing;H,\varnothing)\).

\emph{(1) iff (3).}  By \cref{thm:p03-w03-p03-038-pointed-ef} with \(n=0\),
which applies because \(\tau\) is finite, Duplicator has a winning strategy
in \(G_t(G,\varnothing;H,\varnothing)\) iff
\((G,\varnothing)\equiv_t(H,\varnothing)\), and for \(n=0\) this relation is
agreement on all sentences of quantifier rank at most \(t\), which is (3).

\emph{(3) implies (2)} because a sentence of depth exactly \(t\) has depth at
most \(t\).

\emph{(2) implies (3).}  Let \(z_1,\ldots,z_t\) be distinct variables and set
\[
 \theta_t:=\exists z_1\,\exists z_2\cdots\exists z_t\,(z_1=z_1).
\]
By
\cref{def:p03-quantifier-rank}, \(\operatorname{qr}(\theta_t)=t\), and
\(\theta_t\) is true in every structure with a nonempty domain.  Let
\(\psi\) be a sentence with \(\operatorname{qr}(\psi)\leq t\).  Then
\[
 \operatorname{qr}(\psi\land\theta_t)=\max(\operatorname{qr}(\psi),t)=t,
\]
and in each of \(G\) and \(H\) the sentence \(\psi\land\theta_t\) has the
same truth value as \(\psi\).  By
(2), \(G\entails\psi\land\theta_t\) iff \(H\entails\psi\land\theta_t\), hence
\(G\entails\psi\) iff \(H\entails\psi\).
\end{proof}
\end{annexitem}

\begin{annexitem}{annex:p03-w03-p03-072-infinite-language}{thm:p03-w03-p03-072-infinite-language}{With infinitely many relation symbols the game theorem fails}
\begin{proof}
\emph{Spoiler wins the one-round game.}  Spoiler chooses \(a_\infty\) in
\(H\).  Duplicator must answer some \(a_n\) in \(G\).  Then
\(H\entails U_{n+1}(a_\infty)\) and \(G\not\entails U_{n+1}(a_n)\), so the
pair \((a_n,a_\infty)\) violates the winning condition for the symbol
\(U_{n+1}\).  So every play is won by Spoiler.

\emph{Agreement on depth at most \(1\).}  The domain of \(G\) is contained in
that of \(H\), and each \(U_k\) has the same extension on it in both, so
\(G\) is a substructure of \(H\).  We show by induction on formulas \(\phi\)
with \(\operatorname{qr}(\phi)\leq1\) that for every assignment \(s\) of
elements of \(G\) to the variables,
\(G\entails\phi[s]\) iff \(H\entails\phi[s]\).  Since \(G\) is nonempty, the
case of sentences follows.

For quantifier-free \(\phi\) this holds because \(G\) is a substructure of
\(H\): each atomic formula \(z=z'\) or \(U_k(z)\) has the same truth value
in both under \(s\), and Boolean combinations preserve agreement.  Negation
and conjunction of formulas of rank at most \(1\) preserve the claim, and
a universal formula \(\forall x\,\psi\) is treated as \(\neg\exists x\,\neg\psi\).  The
remaining case is \(\phi=\exists x\,\psi\) with \(\psi\) quantifier-free,
since \(\operatorname{qr}(\phi)\leq1\).  If \(G\entails\phi[s]\), a witness
\(c\in G\) gives \(G\entails\psi[s(x\mapsto c)]\), hence
\(H\entails\psi[s(x\mapsto c)]\) by the quantifier-free case, and
\(H\entails\phi[s]\).  Conversely let \(H\entails\psi[s(x\mapsto c)]\) with
\(c\in H\).  If \(c\in G\), the quantifier-free case gives
\(G\entails\phi[s]\).  If \(c=a_\infty\), let \(m\geq1\) exceed every index
\(k\) of a symbol \(U_k\) occurring in \(\psi\), and choose \(N\geq m\) such
that \(a_N\) differs from \(s(z)\) for each of the finitely many variables
\(z\neq x\) occurring in \(\psi\).  Compare the atomic subformulas of
\(\psi\) under \(s(x\mapsto a_\infty)\) and \(s(x\mapsto a_N)\) in \(H\).
Atoms without \(x\) are unchanged.  \(x=x\) is true under both.  For
\(z\neq x\), \(x=z\) and \(z=x\) are false under both, since
\(s(z)\in G\) while \(a_\infty\notin G\), and \(s(z)\neq a_N\) by the choice
of \(N\).  \(U_k(x)\) with \(k<m\leq N\) is true under both.  So
\(H\entails\psi[s(x\mapsto a_N)]\), and by the quantifier-free case
\(G\entails\psi[s(x\mapsto a_N)]\); hence \(G\entails\phi[s]\).
\end{proof}
\end{annexitem}

\begin{annexitem}{annex:p03-w03-p03-072-transfer}{thm:p03-w03-p03-072-transfer}{Almost sure and complete gives zero--one}
\begin{proof}
\emph{Item 1.}  Let \(B\) be a sentence with \(T\entails B\).  By clause 2 of
\cref{thm:fol-compactness} there are \(A_1,\ldots,A_m\in T\) with
\(A_1,\ldots,A_m\entails B\) (if \(m=0\), \(B\) is valid and
\(\Pr_n[B]=1\)).  Every structure that fails \(B\) fails some \(A_i\), so
\[
  \Pr_n[\neg B]\leq\Pr_n[\neg A_1\lor\cdots\lor\neg A_m]
  \leq\sum_{i=1}^m\Pr_n[\neg A_i]\longrightarrow0 ,
\]
since each \(A_i\) holds almost surely and \(m\) does not depend on \(n\).  So
\(B\) holds almost surely.  Taking for \(B\) the conjunction of finitely many
members of \(T\), its probability tends to \(1\), so it is positive for some
\(n\), and the conjunction has a model with \(n\) elements.  Every finite
subset of \(T\) therefore has a model, and \(T\) has one by
\cref{thm:fol-compactness}.

\emph{Item 2.}  Let \(B\) be any sentence.  As \(T\) is complete,
\(T\entails B\) or \(T\entails\neg B\).  In the first case \(B\) holds almost
surely by item 1; in the second \(\neg B\) holds almost surely, so \(B\)
holds almost never.  If \(B\) holds almost surely and \(T\not\entails B\),
then \(T\entails\neg B\) and \(\Pr_n[B]\) would tend both to \(1\) and to
\(0\), which is impossible.  Hence \(B\) holds almost surely iff
\(T\entails B\), and the almost sure sentences are exactly the consequences
of \(T\).
\end{proof}
\end{annexitem}

\begin{annexitem}{annex:p03-w03-p03-072-closure-strategy}{thm:p03-w03-p03-072-closure-strategy}{Duplicator's closure strategy}
\begin{proof}
Closures, rigidity and density are taken at \(\lambda\) in both graphs, in
the sense of \cref{def:p03-w03-p03-070-weight,def:p03-w03-p03-070-closure}.

\emph{The bounds are admissible.}  We check that \(t_0,\ldots,t_u\) is
admissible for \((\lambda,\lambda,u)\) in the sense of
\cref{def:p03-w03-p03-070-admissible}.  By definition \(t_0=0<t_1=1\), and
\(t_{j+1}\geq t_j\) because \(t_{j+1}\) is a maximum with \(t_j\) among its
terms.  Condition (A1) with \(\mu=\lambda\) says that \(v<e\lambda\) iff
\(v<e\lambda\).  Condition (A2) with \(\kappa=\lambda\) says that
\(t_{j+1}\geq\lfloor(u-j)/\varepsilon_\lambda(t_j)\rfloor+1\) for
\(1\leq j\leq u-1\), which holds because this number is the second term of
the maximum defining \(t_{j+1}\).  (For \(j\geq1\), \(t_j\geq1\), so
\(\varepsilon_\lambda(t_j)\) is defined and positive.)

\emph{The strategy.}  For \(0\leq i\leq u\) and \(s=u-i\) let (I\(_i\)) be the
invariant of item 1 after round \(s\).  We show that (I\(_u\)) holds before
the first round and that, whenever (I\(_i\)) holds with \(i\geq1\), every
Spoiler move in round \(s+1\) has an answer after which (I\(_{i-1}\)) holds.
Duplicator's strategy is to choose such an answer; for infinite graphs the
selection uses the Axiom of Choice.

\emph{Start.}  Let \(Y\) be a set of vertices of \(G\) with
\(1\leq|Y|\leq t_u\).  If \(Y\) were \(\lambda\)-rigid over \(\varnothing\),
then, taking \(Y_0=\varnothing\) in \cref{def:p03-w03-p03-070-weight}, \(Y\)
would be \(\lambda\)-dense over \(\varnothing\):
\(|Y|-\lambda\,e_G(\varnothing,Y)<0\), where \(e_G(\varnothing,Y)\) is the
number of edges of \(G|_Y\).  Then \(G|_Y\) would be a copy of a finite graph
\(F\) with \(|V(F)|-\lambda|E(F)|<0\), which the Nonexistence schema of
\(T^-_\lambda\) excludes.  So \(\varnothing\) has no rigid
\(t_u\)-extension, it is \(t_u\)-closed, and
\(\operatorname{cl}_{t_u}(\varnothing)=\varnothing\) in \(G\), being the
intersection of the \(t_u\)-closed sets containing \(\varnothing\).  The same
holds in \(G'\).  For \(s=0\) the tuples are empty and the empty map is the
required isomorphism.  So (I\(_u\)) holds.

\emph{Step.}  Let \(1\leq i\leq u\), \(s=u-i\), and suppose (I\(_i\)) holds
with tuples \(\bar x=(x_1,\ldots,x_s)\), \(\bar x'=(x'_1,\ldots,x'_s)\) and
isomorphism \(\varphi\).  The hypotheses of
\cref{thm:p03-w03-p03-070-extension} hold with \(\mu=\lambda\): the bounds
are admissible for \((\lambda,\lambda,u)\), \(G\) and \(G'\) are models of
\(T^-_\lambda\), \(s\leq u-i\), and \(\varphi\) is an isomorphism from
\(G|_{\operatorname{cl}_{t_i}(\bar x)}\) onto
\(G'|_{\operatorname{cl}_{t_i}(\bar x')}\) with \(\varphi(x_j)=x'_j\).  If
Spoiler chooses \(y\) in \(G\), the lemma gives \(y'\in G'\) and an
isomorphism \(\varphi_1\) from
\(G|_{\operatorname{cl}_{t_{i-1}}(\bar xy)}\) onto
\(G'|_{\operatorname{cl}_{t_{i-1}}(\bar x'y')}\) with
\(\varphi_1(x_j)=x'_j\) and \(\varphi_1(y)=y'\); Duplicator answers \(y'\).
If Spoiler chooses \(y'\) in \(G'\), the last sentence of the lemma, with the
roles of \(G\) and \(G'\) exchanged and \(\varphi^{-1}\) in place of
\(\varphi\), gives \(y\in G\) and an isomorphism \(\psi\) from
\(G'|_{\operatorname{cl}_{t_{i-1}}(\bar x'y')}\) onto
\(G|_{\operatorname{cl}_{t_{i-1}}(\bar xy)}\) with \(\psi(x'_j)=x_j\) and
\(\psi(y')=y\); Duplicator answers \(y\), and \(\varphi_1=\psi^{-1}\).  In both
cases \(x_{s+1}=y\), \(x'_{s+1}=y'\), \(s+1=u-(i-1)\), and \(\varphi_1\)
witnesses (I\(_{i-1}\)).  This proves item 1.

\emph{Item 2.}  After round \(u\), (I\(_0\)) holds with \(t_0=0\).  Every set
is \(0\)-closed, because a rigid \(0\)-extension would have at least one and
at most \(0\) elements; so \(\operatorname{cl}_0(X)=X\) for every \(X\).
Hence \(\varphi\) is an isomorphism from \(G|_{\{x_1,\ldots,x_u\}}\) onto
\(G'|_{\{x'_1,\ldots,x'_u\}}\) with \(\varphi(x_j)=x'_j\).  As \(\varphi\) is
a well-defined injective map, \(x_j=x_k\) iff \(x'_j=x'_k\); as it preserves
adjacency in both directions, \(x_j\sim x_k\) iff \(x'_j\sim x'_k\).  This is
the winning condition of \(\mathrm{EHR}(G,G';u)\).

\emph{Item 3.}  Models of \(T^-_\lambda\) are nonempty: the rooted graph with
no roots and one nonroot \(y\) has \(v=1\), \(e=0\), and its only nonempty
subset of nonroots, \(\{y\}\), has weight \(1>0\), so it is
\(\lambda\)-safe; its Generic Extension sentence with \(t=0\) asserts that a
vertex exists.  Let \(G,G'\) be models of \(T^-_\lambda\).  By item 2
Duplicator wins \(\mathrm{EHR}(G,G';u)\) for every \(u\geq1\).  The language
\(\{\sim\}\) has one relation symbol, so by
\cref{thm:p03-w03-p03-072-game-theorem}, (1) to (3), \(G\) and \(G'\) agree
on all sentences of depth at most \(u\), for every \(u\); every sentence has
a finite depth, so they are elementarily equivalent.  If \(T^-_\lambda\) is
consistent and \(B\) is a sentence with \(T^-_\lambda\not\entails B\) and
\(T^-_\lambda\not\entails\neg B\), then \(T^-_\lambda\cup\{\neg B\}\) and
\(T^-_\lambda\cup\{B\}\) have models, and these two models of
\(T^-_\lambda\) are not elementarily equivalent, which is impossible.  So
\(T^-_\lambda\) is complete.
\end{proof}
\end{annexitem}

\begin{annexitem}{annex:p03-w03-p03-072-sparse-zero-one}{thm:p03-w03-p03-072-sparse-zero-one}{The zero--one law for \(n^{-\alpha}\) by the game}
\begin{proof}
For each \(n\), \(G(n,n^{-\alpha})\) (\cref{def:p03-w03-p03-065-gnp}) is a
probability measure on the graphs with vertex set \(\{1,\ldots,n\}\), that
is, on the \(\{\sim\}\)-structures with that domain in which \(\sim\) is
irreflexive and symmetric; so it is a random \(\{\sim\}\)-structure in the
sense of \cref{def:p03-w03-p03-072-random-structure}, and almost sure truth
there is that of \cref{def:p03-w03-p03-065-almost-sure}.  Since
\(\alpha\) is irrational, \cref{thm:p03-w03-p03-070-as-holds} gives
\((\mathrm{AS}_\alpha)\): every member of \(T_\alpha\) holds almost surely.
By \cref{thm:p03-w03-p03-072-transfer}, item 1, \(T_\alpha\) is consistent.
By \cref{thm:p03-w03-p03-072-closure-strategy}, item 3, with
\(\lambda=\alpha\), it is complete.  By
\cref{thm:p03-w03-p03-072-transfer}, item 2, every sentence \(B\) holds
almost surely or almost never, so
\(\lim_{n\to\infty}\Pr[G(n,n^{-\alpha})\entails B]\) exists and lies in
\(\{0,1\}\), and it equals \(1\) iff \(T_\alpha\entails B\).
\end{proof}
\end{annexitem}
